\section{Фундаментальные узлы компьютерных систем и двоичные функции.}

Реализация компьютерных систем базируется на том, что любая булева функция может быть воплощена в виде электронной схемы. В основе этого лежит связь между «малой» арифметикой (логические операции над битами) и «большой» арифметикой (операции над машинными словами).

\subsection*{Логические вентили (Базовые элементы)}
Каждая элементарная булева функция реализуется физическим устройством, называемым \textbf{вентилем}:
\begin{itemize}
    \item \textbf{Инвертор (NOT):} реализует $\neg x$.
    \item \textbf{Конъюнктор (AND):} реализует $x \land y$.
    \item \textbf{Дизъюнктор (OR):} реализует $x \lor y$.
    \item \textbf{Исключающее ИЛИ (XOR):} реализует $x \oplus y$ (основа сумматоров).
\end{itemize}

\subsection*{Комбинационные узлы}
Узлы, значение на выходе которых зависит только от текущих значений на входе (реализация булевых функций «в чистом виде»).

\paragraph{1. Сумматор (Основа арифметики)}
Служит для сложения двоичных чисел. Строится на основе разложения суммы по разрядам.
\begin{itemize}
    \item \textbf{Полусумматор (Half Adder):} Складывает два бита $x$ и $y$.
    \begin{itemize}
        \item Сумма: $S = x \oplus y$
        \item Перенос: $C = x \land y$
    \end{itemize}
    \item \textbf{Полный сумматор (Full Adder):} Складывает два бита $x, y$ и перенос из предыдущего разряда $C_{in}$.
    \begin{itemize}
        \item $S = x \oplus y \oplus C_{in}$
        \item $C_{out} = (x \land y) \lor (C_{in} \land (x \oplus y))$
    \end{itemize}
\end{itemize}

\paragraph{2. Дешифратор (Decoder)}
Узел, который преобразует двоичный код в сигнал на одной из выходных линий.
\begin{itemize}
    \item Если на вход подано $n$ бит, то выходов будет $2^n$.
    \item Каждый выход соответствует определенной \textbf{элементарной конъюнкции} (минтерму) входных переменных.
\end{itemize}

\paragraph{3. Мультиплексор (Multiplexer)}
Узел, который выбирает один из нескольких информационных входов и передает его на единственный выход. Выбор осуществляется на основе адресного кода.
\begin{itemize}
    \item Реализует \textbf{разложение Шеннона}: адресные линии выступают в роли переменных разложения, выбирая нужную подфункцию.
\end{itemize}

\subsection*{Узлы с памятью (Последовательностные)}
Для реализации полноценных вычислений необходимы элементы, хранящие состояние (память).

\paragraph{RS-триггер}
Простейший элемент памяти, строящийся на обратной связи (например, на двух элементах ИЛИ-НЕ / Стрелках Пирса).
\begin{itemize}
    \item Позволяет хранить 1 бит информации.
    \item В отличие от комбинационных схем, здесь значение функции зависит не только от входов $R$ (Reset) и $S$ (Set), но и от предыдущего значения выхода $Q$.
\end{itemize}

\subsection*{Синтез схем}
Согласно теореме о функциональной полноте (билет 32), используя только вентили И, ИЛИ, НЕ (или только штрих Шеффера), можно построить узел любой сложности.
\begin{itemize}
    \item \textbf{Прямой синтез:} По таблице истинности строится СДНФ, которая напрямую переводится в двухуровневую логическую схему (слой конъюнкторов $\to$ слой дизъюнкторов).
    \item \textbf{Минимизация:} Использование МДНФ (билет 28) позволяет сократить количество вентилей и входов, что критично для стоимости и скорости работы процессора.
\end{itemize}